\documentclass[tikz,border=20pt]{standalone}
\usepackage{tikz}
\usetikzlibrary{arrows.meta, positioning, fit, backgrounds}

\begin{document}
\begin{tikzpicture}[
    font=\small\sffamily,
    % Box styling with balanced, clean padding
    box/.style={
        draw=black!85,
        thick,
        fill=white,
        rounded corners=3pt,
        align=center,
        inner sep=9pt
    },
    flowArrow/.style={
        -{Stealth[length=7pt, width=5pt]},
        very thick,
        draw=black!85
    }
]

    % =========================================================================
    % LAYER 1: USER INTERFACE LAYER (y = 0)
    % =========================================================================
    \node[box, fill=blue!8, text width=5.6cm] (cli) at (-4.0, 0) {
        \textbf{\large CLI Frontend (\texttt{synth})}\\[3pt]
        \scriptsize Parses CLI args (\texttt{--freq}, \texttt{--wave}, etc.),\\validates ranges, triggers audio rendering.
    };

    \node[box, fill=blue!8, text width=5.6cm] (tui) at (4.0, 0) {
        \textbf{\large Interactive TUI (\texttt{synth-tui})}\\[3pt]
        \scriptsize FTXUI interactive GUI with real-time waveform canvas,\\ADSR envelope graph, and parameter sliders.
    };

    \begin{scope}[on background layer]
        \node[fill=gray!8, draw=black!45, dashed, thick, rounded corners=6pt,
              fit=(cli) (tui), inner ysep=20pt, inner xsep=16pt] (layer1) {};
        % Title INSIDE the dotted box at top-left
        \node[anchor=north west, font=\bfseries\footnotesize, text=black!85]
              at ([xshift=10pt, yshift=-7pt]layer1.north west) {1. User Interface Layer (Executables)};
    \end{scope}

    % =========================================================================
    % LAYER 2: CORE SYNTHESIZER ENGINE (y = -5.0 to -8.6)
    % =========================================================================
    \node[box, fill=indigo!8, text width=5.6cm] (settings) at (-4.0, -5.0) {
        \textbf{\large Configuration / Settings}\\[3pt]
        \scriptsize \texttt{SynthSettings} struct \& \texttt{Waveform} enum\\Frequency, duration, sample rate, ADSR.
    };

    \node[box, fill=indigo!8, text width=5.6cm] (pipeline) at (4.0, -5.0) {
        \textbf{\large Rendering Pipeline (\texttt{render})}\\[3pt]
        \scriptsize Computes output audio sample buffer:\\[2pt]
        $\text{sample}[i] = \text{osc} \cdot \text{env} \cdot \text{amp}$
    };

    % DSP Core: Correct subsystem descriptions
    \node[box, fill=white, text width=12.4cm] (dspCore) at (0, -8.6) {
        \textbf{\large DSP Core Sound Generation Subsystems (\texttt{synth\_core})}\\[5pt]
        \begin{tabular}{p{5.6cm}@{\hspace{0.8cm}}p{5.6cm}}
            \centering\textbf{\normalsize Oscillator Hierarchy} & \centering\textbf{\normalsize ADSR Envelope Generator}\\[2pt]
            \centering\scriptsize Sine, Square, Sawtooth periodic generators & \centering\scriptsize 4-Stage piecewise linear modulator
        \end{tabular}
    };

    \begin{scope}[on background layer]
        \node[fill=blue!4, draw=blue!50!black, dashed, thick, rounded corners=6pt,
              fit=(settings) (pipeline) (dspCore), inner ysep=24pt, inner xsep=16pt] (layer2) {};
        % Title INSIDE the dotted box — two compact lines staying well clear of dataflow arrows
        \node[anchor=north west, font=\bfseries\footnotesize, text=black!85]
              at ([xshift=10pt, yshift=-7pt]layer2.north west) {\begin{tabular}{@{}l@{}}\textbf{2. Core Synthesizer Engine}\\\scriptsize\texttt{(synth\_core Static Library)}\end{tabular}};
    \end{scope}

    % =========================================================================
    % LAYER 3: STORAGE & FILE OUTPUT (y = -12.6)
    % =========================================================================
    \node[box, fill=green!8, text width=6.6cm] (wavWriter) at (-2.0, -12.6) {
        \textbf{\large WAV Serializer (\texttt{WavWriter})}\\[3pt]
        \scriptsize 16-bit PCM quantization, Little-Endian bit encoding,\\44-byte standard RIFF/WAVE header assembly.
    };

    \node[box, fill=green!18, text width=4.0cm] (wavFile) at (5.0, -12.6) {
        \textbf{\large Output File}\\[3pt]
        \scriptsize \texttt{tone.wav}\\(Standard 16-bit PCM)
    };

    \begin{scope}[on background layer]
        \node[fill=green!4, draw=green!60!black, dashed, thick, rounded corners=6pt,
              fit=(wavWriter) (wavFile), inner ysep=24pt, inner xsep=16pt] (layer3) {};
        % Title INSIDE the dotted box — two compact lines staying well clear of dataflow arrows
        \node[anchor=north west, font=\bfseries\footnotesize, text=black!85]
              at ([xshift=10pt, yshift=-7pt]layer3.north west) {\begin{tabular}{@{}l@{}}\textbf{3. Storage and Output Layer}\\\scriptsize\texttt{(RIFF/WAVE Audio File)}\end{tabular}};
    \end{scope}

    % =========================================================================
    % CLEAN DATAFLOW ARROWS (Zero overlap with text, boxes, or layer titles)
    % =========================================================================
    % UI Bus Bar: Clean horizontal bus in open space between Layer 1 and Layer 2 (at y = -2.2)
    \draw[very thick, draw=black!85] (cli.south) -- (-4.0, -2.2);
    \draw[very thick, draw=black!85] (tui.south) -- (4.0, -2.2);
    \draw[very thick, draw=black!85] (-4.0, -2.2) -- (4.0, -2.2);

    % Single clean arrow dropping from the bus bar into Configuration / Settings
    % Drops vertically at x = -1.6, which is > 1.6 cm clear to the right of Layer 2 title
    \draw[flowArrow] (-1.6, -2.2) -- (-1.6, -2.2 |- settings.north);

    % Settings -> Rendering Pipeline (straight horizontal across 2.4 cm gap)
    \draw[flowArrow] (settings.east) -- (pipeline.west);

    % Rendering Pipeline -> DSP Core (straight vertical drop at x = 4.0)
    \draw[flowArrow] (pipeline.south) -- (pipeline.south |- dspCore.north);

    % DSP Core -> WavWriter (clean straight vertical drop at x = -0.6)
    % Drops at x = -0.6, which is > 2.2 cm clear to the right of Layer 3 title
    % Label sits comfortably in the 2.4 cm vertical gap between Layer 2 and Layer 3
    \draw[flowArrow] ([xshift=-0.6cm]dspCore.south) -- node[midway, fill=white, inner sep=3pt, font=\scriptsize\sffamily] {Buffer: \texttt{std::vector<double>}} ([xshift=1.4cm]wavWriter.north);

    % WavWriter -> Output File (clean horizontal arrow with label neatly placed above)
    \draw[flowArrow] (wavWriter.east) -- node[midway, above=3pt, font=\scriptsize\sffamily] {Binary Stream} (wavFile.west);

\end{tikzpicture}
\end{document}
